%! TEX root = ../am2-20-21.tex

\chapter{Recapitulare examen}

\indent\indent\textbf{MODEL 1}

1. Folosind eventual formula Green-Riemann, calculați
$ \ds\int_\gamma 3xy^2 dx + 2x^2y dy $, unde $ \gamma $
este triunghiul $ \Delta ABC $, cu $ A(1, 1), B(3, 3), C(3, 1) $.
Reprezentare grafică.

\vspace{0.3cm}

2. Calculați aria suprafeței $ \Sigma: 2y = 3x^2 + 3z^2, y \leq 2 $.
Reprezentare grafică.

\vspace{0.3cm}

3. Calculați fluxul cîmpului vectorial $ \vec{V} = 3x \vec{i} + (x^2 + z^2)\vec{j} - 2z\vec{k} $,
prin suprafața exterioară a emisferei $ \Sigma: x^2 + y^2 + z^2 = 4, z \leq 0 $,
orientată cu normala spre exterior. Reprezentare grafică.

\vspace{0.3cm}

4. Calculați circulația cîmpului vectorial $ \vec{V} = x\vec{i} + 2y\vec{j} + z^2 \vec{k} $
de-a lungul curbei de intersecție între suprafețele $ \Sigma_1: z = x^2 + y^2 $
și $ \Sigma_2: z = 4 $. Reprezentare grafică.

\vspace{0.3cm}

5. Să se dezvolte în serie Fourier funcția $ f(x) = x^3, x \in (-\pi, \pi) $.

\vspace{2cm}

\textbf{MODEL 2}

1. Folosind eventual formula Green-Riemann, calculați $ \ds\int_\gamma (x^2 + y^2) dx + (x^2 - y^2) dy $,
unde $ \gamma: x^2 - 4x + y^2 + 2y = 0, y \geq 0 $. Reprezentare grafică.

\vspace{0.3cm}

2. Calculați aria suprafeței $ \Sigma: x^2 = 3y^2 + 3z^2, 0 \leq x \leq 2 $. Reprezentare grafică.

\vspace{0.3cm}

3. Calculați fluxul cîmpului vectorial $ \vec{V} = y^2 \vec{i} + 2x^2 \vec{j} + 3z^2 \vec{k} $
prin suprafața exterioară a paraboloidului $ \Sigma: z = 2x^2 + 2y^2, z \leq 2 $,
orientată cu normala spre exterior. Reprezentare grafică.

\vspace{0.3cm}

4. Calculați circulația cîmpului vectorial $ \vec{V} = (x + y) \vec{i} + (x + z) \vec{j} + (y + z) \vec{k} $
de-a lungul curbei de intersecție între suprafețele $ \Sigma_1: x^2 + y^2 + z^2 = 4 $ și
$ \Sigma_2: y = 1 $. Reprezentare grafică.

\vspace{0.3cm}

5. Să se dezvolte în serie Fourier funcția $ f(x) = |x^3|, x \in (-\pi, \pi) $.

\vspace{2cm}

\textbf{MODEL 3}

1. Folosind eventual formula Green-Riemann, calculați $ \ds\int_\gamma x^3 dx + (x^2 + y) dy $,
unde $ \gamma: 3x^2 + 2y^2 = 5 $. Reprezentare grafică.

\vspace{0.3cm}

2. Calculați aria suprafeței $ \Sigma: y^2 = 2x^2 + 2z^2, 0 \leq y \leq 3 $. Reprezentare grafică.

\vspace{0.3cm}

3. Calculați fluxul cîmpului vectorial $ \vec{V} = z^2 \vec{i} + 2y^2 \vec{j} + x^2 \vec{k} $
prin suprafața exterioară a cilindrului $ \Sigma: x^2 + y^2 = 4, z \in [1, 3] $, orientat
cu normala spre exterior. Reprezentare grafică.

\vspace{0.3cm}

4. Calculați circulația cîmpului vectorial $ \vec{V} = (x, 2y, y-z) $ de-a lungul
curbei de intersecție între suprafețele $ \Sigma_1: z = 3x^2 + 3y^2, \Sigma_2: z = 3 $.
Reprezentare grafică.

\vspace{0.3cm}

5. Să se dezvolte în serie Fourier funcția $ f(x) = 1 - x, x \in (0, 1) $.